\section{Introduction}
When two agents check each other's work, they may begin approving answers they
have not properly checked. Fixing the setup that allowed this behavior is only
the first step. Do the agents return to useful checking, or does the earlier
pattern persist?

We study this recovery question in a setting where agents must exchange complete
tool logs before approving each other's work. Their messages are too short to
carry those logs. Both completed pairs repeatedly approved without the required
evidence. This gives us an observed failure to test, though it does not establish
a negotiated agreement or show that one agent caused the other's behavior.

Building on \citet{shi2026emergent}, the agents solve tasks with tools, exchange
messages, receive feedback, and keep their own conversations. We then use one
pair's saved conversations to compare three changes: providing complete logs,
reminding agents of the original rule, and starting fresh conversations.
These tests examine one way to repair a flawed setup, not every route to
agent collusion.

Recovery must mean more than fewer approvals. An agent that rejects everything
avoids approving faulty work but also stops being useful. We therefore measure
whether agents reject incorrect work \emph{and} approve correct work. We test
both the next review and four consecutive reviews after a change is introduced
once, using the same tasks across conditions.
